\documentclass[10pt,a4paper]{article}
\usepackage{fontspec,amsmath,amssymb,graphicx}
\usepackage[margin=18mm]{geometry}
\setmainfont{Times New Roman}
\setlength{\parindent}{0pt}
\setlength{\parskip}{5pt}
\allowdisplaybreaks
\sloppy
\emergencystretch=2em
\begin{document}
\section*{VD\_3\_5\_2\_02 — Lời giải bổ sung}
\subsection*{Kiến thức và trình tự giải}
\textbf{Bước 1 — Chọn năng lượng theo toán tử.} Với sóng dùng $E=\frac12\int(u_t^2+a^2u_x^2)dx$. Với Klein–Gordon có thêm $u$, thêm $\frac12\int u^2dx$. Với dầm $u_{tt}=-a^2u_{xxxx}$ dùng $E=\frac12\int(u_t^2+a^2u_{xx}^2)dx$, không thay $u_{xx}$ bằng $u_x$.

\textbf{Bước 2 — Đạo hàm và tích phân từng phần.} Với $u_{tt}-a^2u_{xx}+\gamma u_t+\kappa u=F$, năng lượng có hạng $\kappa u^2$ cho


\[
E'=a^2[u_tu_x]_0^L-\gamma\int u_t^2dx+\int Fu_tdx.
\]


Với dầm, tích phân hai lần cho $E'=a^2[u_{tx}u_{xx}-u_tu_{xxx}]_0^L$. Viết riêng từng hạng biên rồi dùng điều kiện biên và đạo hàm theo thời gian của điều kiện đó để cho chúng bằng không.

\textbf{Bước 3 — Suy hệ quả.} Không cản, không nguồn thì $E$ bảo toàn. Có cản $\gamma>0$ thì $E'\le0$; nếu $E$ không có hệ số $1/2$, $E'=-2\gamma\int u_t^2\ge-2\gamma E$, suy $e^{-2\gamma t}E(0)\le E(t)\le E(0)$. Với hiệu hai nghiệm cùng dữ liệu, $E(0)=0$ nên có duy nhất.

\textbf{Bước 4 — Tách biến nếu phải tìm nghiệm cụ thể.} Sóng có cản $\gamma=1$ cho $T''+T'+a^2\lambda T=0$. Đặt $T=e^{-t/2}Y$, thu $Y''+(a^2\lambda-1/4)Y=0$, $Y(0)=A$, $Y'(0)=B+A/2$. Phân biệt dao động, nghiệm kép và hàm hyperbolic. Dầm cho $X''''=k^4X$, $T''+a^2k^4T=0$; dùng đủ bốn điều kiện biên để tìm phương trình của $kL$.

\textbf{Bước 5 — Kiểm tra hệ số và khẳng định của đề.} Năng lượng đầu chứa cả $g^2$ và các đạo hàm của $f$. Thử $f=0$, $g\ne0$ để phát hiện công thức bỏ vận tốc. Với dầm, nghiệm $k=0$ phải giải trực tiếp bằng đa thức bậc ba; nghiệm $0$ của phương trình lượng giác–hyperbolic có thể chỉ là nghiệm giả do cơ sở suy biến.

\textbf{Chi tiết phép triệt tiêu tích phân.} Đạo hàm động năng cho $\int u_tu_{tt}$; dùng PDE thay $u_{tt}$. Với sóng, $\int u_tu_{xx}=[u_tu_x]-\int u_{tx}u_x$, và tích phân cuối triệt tiêu đạo hàm thế năng $\int u_xu_{xt}$. Với dầm,


\[
\int u_tu_{xxxx}=[u_tu_{xxx}-u_{tx}u_{xx}]+\int u_{txx}u_{xx}.
\]


Tích phân cuối triệt tiêu đạo hàm $\frac12\int u_{xx}^2$ sau khi nhân hệ số và giữ đúng dấu PDE. Viết hạng biên trước, sau đó mới dùng điều kiện của từng đầu.

\textbf{Bước 1 — Tách biến và chọn bài toán trị riêng theo biên.} Đặt $u=X(x)T(t)$ trong $u_{tt}=a^2u_{xx}$, suy $X''+\lambda X=0$, $T''+a^2\lambda T=0$. Tích phân từng phần cho $\lambda\int X^2=\int (X')^2$ khi hạng biên bằng không; do đó $\lambda\ge0$. Giải cả trường hợp $\lambda=0$ trước khi xét $\lambda=k^2>0$.

\textbf{Bước 2 — Tìm hệ hàm riêng.} Trên $(0,L)$: Dirichlet hai đầu dùng $\sin(n\pi x/L)$, $n\ge1$; Neumann hai đầu dùng $1$ và $\cos(n\pi x/L)$; Dirichlet–Neumann dùng $\sin((n+1/2)\pi x/L)$, $n\ge0$; Neumann–Dirichlet dùng cos của các tần số nửa nguyên ấy. Chuẩn bình phương của mode khác không là $L/2$, của mode hằng là $L$.

\textbf{Bước 3 — Chiếu cả hai dữ liệu.} Với mode khác không đặt


\[
A_n=\frac2L\int_0^L fX_n\,dx,\quad B_n=\frac2L\int_0^L gX_n\,dx,\quad \omega_n=ak_n.
\]


Tích phân từng khoảng nếu dữ liệu chia đoạn; dùng tích phân từng phần hoặc công thức tích–tổng khi là đa thức/lượng giác. Hệ số $B_n$ là hệ số của vận tốc, nên hệ số đứng trước $\sin\omega_nt$ phải là $B_n/\omega_n$.

\textbf{Bước 4 — Ghép nghiệm và xét mode hằng.}


\[
u=\sum_n\left(A_n\cos\omega_nt+\frac{B_n}{\omega_n}\sin\omega_nt\right)X_n.
\]


Với Neumann hai đầu, thêm $A_0+B_0t$, trong đó $A_0=L^{-1}\int f$, $B_0=L^{-1}\int g$. Với tần số bị thay đổi bởi hạng $u$ hoặc cản $u_t$, giải lại ODE của từng mode; không dùng nguyên tần số $ak_n$.

\textbf{Bước 5 — Kiểm tra biên, dữ liệu và hội tụ.} Mỗi mode thỏa biên; tại $t=0$ hai chuỗi cho khai triển của $f,g$. Nếu dữ liệu không tương thích ở góc, giải theo nghĩa năng lượng hoặc từng miền và nêu đường mất tính trơn; không đạo hàm một chuỗi yếu tùy ý. Chu kỳ được suy từ tất cả mode có hệ số khác không, kể cả mode hằng; hai tần số có tỷ số vô tỷ ngăn một chu kỳ chung.

\textbf{Chi tiết tính tích phân hệ số.} Với $k=n\pi/L$, hai nguyên hàm thường dùng là


\[
\int x\sin kx\,dx=-x\cos kx/k+\sin kx/k^2,\qquad
\int x\cos kx\,dx=x\sin kx/k+\cos kx/k^2.
\]


Với đa thức $P$, ghi cả hai lần tích phân:


\[
\int_0^L P\sin kx\,dx
=\left[-P\cos kx/k\right]_0^L+
\left[P'\sin kx/k^2\right]_0^L-
\frac1{k^2}\int_0^L P''\sin kx\,dx.
\]


Chỉ bỏ hạng biên sau khi thế $0,L$ và điều kiện thực sự của $P$. Nếu dữ liệu chia đoạn, áp dụng công thức trên từng đoạn, rồi cộng các hạng tại điểm chia. Cuối cùng nhân nghịch đảo chuẩn của hàm riêng; nếu hệ số là của vận tốc, chia thêm tần số khi ghép nghiệm. Với mẫu số bằng $0$ trong tích–tổng, quay lại tích phân gốc để tính riêng.
\subsection*{Lời giải áp dụng cho bài}
\textbf{Bước 1. Năng lượng đúng của dầm.} $u_{tt}=-a^2u_{xxxx}$, hai đầu ngàm $u=u_x=0$. Nhân với $u_t$ và tích phân hai lần:
\[E=\tfrac12\int_0^L(u_t^2+a^2u_{xx}^2)dx,\qquad E'=a^2[u_{tx}u_{xx}-u_tu_{xxx}]_0^L=0.\]
Do $u_t=u_{tx}=0$ ở hai đầu. Năng lượng có $u_x^2$ trong đề không phù hợp toán tử cấp bốn; cần sửa thành $u_{xx}^2$. Với hiệu hai nghiệm, $E=0$ suy $u_t=u_{xx}=0$; hàm affine còn lại bị biên ngàm và dữ liệu đầu buộc bằng $0$, nên duy nhất.

\textbf{Bước 2. Trị riêng dương và trực giao.} $X''''=\lambda X$, $X=X'=0$ ở hai đầu. Tích phân hai lần cho $\lambda\int X^2=\int(X'')^2$. Nếu $\lambda=0$, $X$ affine rồi biên loại hết; vậy $\lambda>0$. Trừ tích phân hai phương trình cho $(\lambda_m-\lambda_n)\int X_mX_n=0$ vì hạng biên $[X_m'''X_n-X_m''X_n'-X_n'''X_m+X_n''X_m']$ bằng $0$.

\textbf{Bước 3. Giải đủ bốn điều kiện.} Đặt $\lambda=k^4$, $\beta=kL$. Hai điều kiện trái cho $X=A(\cosh kx-\cos kx)+B(\sinh kx-\sin kx)$. Hai điều kiện phải tạo định thức $2(1-\cosh\beta\cos\beta)=0$, nên $\cosh\beta\cos\beta=1$. Chỉ lấy nghiệm dương; $\beta=0$ là nghiệm giả của cơ sở suy biến. Nghiệm đầu $\beta_1\approx4.730040745$. Chọn
\[X_n=\cosh k_nx-\cos k_nx-\sigma_n(\sinh k_nx-\sin k_nx),\quad
\sigma_n=\frac{\cosh\beta_n-\cos\beta_n}{\sinh\beta_n-\sin\beta_n}.\]
Phương trình đặc trưng bảo đảm cả $X_n(L)=X_n'(L)=0$.

\textbf{Bước 4. Chiếu dữ liệu và ghép.} $A_n=\int fX_n/\int X_n^2$, $B_n=\int gX_n/\int X_n^2$, tần số $\omega_n=ak_n^2$. Nghiệm $u=\sum(A_n\cos\omega_nt+B_n\sin\omega_nt/\omega_n)X_n$. Kiểm tra bốn biên bằng hàm riêng, hai dữ liệu bằng hệ số; không dùng tần số $ak_n$ của phương trình sóng cấp hai. Với dữ liệu đủ trơn và tương thích thu nghiệm cổ điển; với dữ liệu năng lượng hiểu chuỗi trong không gian thích hợp.

\textbf{Chi tiết hai tích phân năng lượng và định thức.} Nhân PDE với $u_t$:
\[\begin{aligned}
\int_0^Lu_tu_{xxxx}dx
&=[u_tu_{xxx}]_0^L-\int_0^Lu_{tx}u_{xxx}dx\\
&=[u_tu_{xxx}-u_{tx}u_{xx}]_0^L+\int_0^Lu_{txx}u_{xx}dx.
\end{aligned}\]
Hạng tích phân cuối bằng $\frac12(d/dt)\int u_{xx}^2$; chuyển sang cùng động năng cho $E'$ đã viết.

Đặt $C=\cosh\beta,c=\cos\beta,S=\sinh\beta,s=\sin\beta$. Hai điều kiện phải, sau chia đạo hàm cho $k$, là
\[\begin{pmatrix}C-c&S-s\\S+s&C-c\end{pmatrix}\binom AB=\binom00.\]
Định thức $(C-c)^2-(S-s)(S+s)=C^2-2Cc+c^2-S^2+s^2=2-2Cc$, vì $C^2-S^2=c^2+s^2=1$. Dòng đầu cho $B/A=-(C-c)/(S-s)$; dòng hai đúng khi định thức bằng $0$. Đây là nguồn gốc của $\sigma_n$.\end{document}